% =============================================================================
% SECCIÓN: EXPERIENCIAS PROFESIONALES — SECTOR PÚBLICO
% =============================================================================
% Uso: activar en index.tex descomentando:
%   % =============================================================================
% SELECTOR: EXPERIENCIAS PROFESIONALES — SECTOR PÚBLICO
% =============================================================================
% Usado por: profiles/sector-publico.tex
%
% Este archivo solo SELECCIONA y ORDENA entradas. El contenido de cada
% experiencia vive en entradas/ (un archivo por redacción); edítalo allí.
% Las entradas sin sufijo son la redacción base compartida; las --sector-publico
% son redacciones adaptadas a este perfil.
% =============================================================================

\sectionTitle{Experiencias Profesionales}{\faSuitcase}

\begin{experiences}

\entradaExperiencia{2026-05_onpe-capacitador}
\emptySeparator

\entradaExperiencia{2026-01_onpe-coordinador}
\emptySeparator

\entradaExperiencia{2025-10_cau-metodologia-ia--sector-publico}
\emptySeparator

\entradaExperiencia{2025-08_inei-censista}
\emptySeparator

\entradaExperiencia{2025-01_cau-metodologia--sector-publico}
\emptySeparator

\entradaExperiencia{2024-12_sofia-coordinador--sector-publico}
\emptySeparator

% Inactiva — activar si el puesto valora docencia en proyectos de inversión:
%\entradaExperiencia{2024-04_unah-gestion-proyectos--sector-publico}
%\emptySeparator

\entradaExperiencia{2022-04_dra-analista-datos--sector-publico}

\end{experiences}

% =============================================================================

\sectionTitle{Experiencias Profesionales}{\faSuitcase}

\begin{experiences}

\experience
{Ahora} {Docente de Metodología de la Investigación}{Corporación Educativa Universitaria}{Ayacucho}
{Enero 2026} {
	Formación académica en metodología de investigación científica con énfasis en rigor metodológico, ética investigativa y redacción de documentos técnicos para el sector público.
	\begin{itemize}
		\item Enseñanza de diseño de investigación cuantitativa, cualitativa y mixta aplicada a ciencias sociales
		\item Talleres de IA aplicada a la investigación: uso ético de ChatGPT, búsqueda académica y gestión bibliográfica
		\item Formación en normas APA 7, prevención de plagio y redacción de informes técnicos y científicos
		\item Sesiones prácticas de análisis de datos con R/SPSS y elaboración de marcos teóricos
		\item Evaluación formativa mediante rúbricas y retroalimentación constructiva continua
	\end{itemize}
}
{\footnotesize{\emph{Herramientas:} Google Classroom, Turnitin, Zotero, Mendeley, R/SPSS, modalidad híbrida}}
\emptySeparator

\experience
{Abril 2026} {Coordinador de Centro Poblado}{Oficina Nacional de Procesos Electorales (ONPE) — ODPE Parinacochas}{Coracora, Ayacucho}
{Enero 2026} {
	Coordinación operativa de mesas de sufragio en el marco de las Elecciones Generales 2026, con responsabilidad directa sobre la gestión de personal, documentación oficial y supervisión del proceso electoral en la ODPE Parinacochas.
	\begin{itemize}
		\item Supervisión del proceso de instalación, desarrollo y cierre de mesas de sufragio bajo protocolos ONPE
		\item Manejo de documentación oficial: actas electorales y cadena de custodia de materiales sensibles
		\item Coordinación con miembros de mesa, personeros y autoridades locales para garantizar el proceso
		\item Gestión de incidencias en campo y reporte a la ODPE según procedimientos institucionales establecidos
		\item Trabajo bajo contrato de locación de servicios con cumplimiento de metas operativas institucionales
	\end{itemize}
}
{\footnotesize{\emph{Entidad:} ONPE — ODPE Parinacochas | Período: enero–abril 2026 | Proceso: Elecciones Generales 2026}}
\emptySeparator

\experience
{Octubre 2025} {Censista — Recolección de Datos Estadísticos}{Instituto Nacional de Estadística e Informática (INEI)}{Ayacucho}
{Agosto 2025} {
	Recolección, validación y reporte de datos socioeconómicos en campo para la Oficina Departamental de Estadística e Informática de Ayacucho (ODEI), en el marco de operaciones estadísticas del Estado.
	\begin{itemize}
		\item Levantamiento de información socioeconómica de hogares según metodología oficial del INEI
		\item Aplicación de instrumentos estructurados con control de calidad y consistencia de datos en campo
		\item Manejo de dispositivos electrónicos y sistemas de registro digital para captura de información
		\item Coordinación con supervisores y equipos técnicos para validación de datos y cobertura de metas
		\item Cumplimiento estricto de protocolos institucionales de confidencialidad y ética estadística
	\end{itemize}
}
{\footnotesize{\emph{Entidad:} INEI — ODEI Ayacucho | Modalidad: Locación de Servicios | Período: agosto–octubre 2025}}
\emptySeparator

\experience
{Julio 2024} {Docente de Gestión de Proyectos II}{Universidad Nacional Autónoma de Huanta}{Huanta}
{Abril 2024} {
	Docencia en gestión y evaluación de proyectos de inversión pública con enfoque cuantitativo y análisis de viabilidad bajo estándares del sector público peruano.
	\begin{itemize}
		\item Enseñanza de análisis costo-beneficio, VAN y TIR aplicados a proyectos de inversión pública
		\item Evaluación económica y financiera de proyectos bajo lineamientos del Sistema Nacional de Inversión Pública
		\item Elaboración de reportes técnicos, fichas de proyecto y presentaciones ejecutivas con indicadores de desempeño
		\item Uso de herramientas de gestión: cronogramas, presupuestos, indicadores y marcos lógicos
		\item Retroalimentación personalizada basada en competencias y rúbricas analíticas
	\end{itemize}
}
{\footnotesize{\emph{Tecnologías utilizadas:} Excel avanzado, MS Project, Trello, Google Workspace, LaTeX}}
\emptySeparator

\experience
{Marzo 2025} {Docente de Metodología de la Investigación}{Corporación Educativa Universitaria}{Ayacucho}
{Enero 2025} {
	Formación en fundamentos metodológicos para el diseño de proyectos de investigación aplicada, con énfasis en ciencias sociales y políticas públicas.
	\begin{itemize}
		\item Enseñanza de epistemología, tipos de investigación y enfoques metodológicos cuantitativos y cualitativos
		\item Diseño de investigaciones aplicadas a problemáticas sociales, económicas y de gestión pública
		\item Elaboración, validación y confiabilización de instrumentos de recolección de datos
		\item Técnicas de muestreo probabilístico y no probabilístico para estudios sociales
		\item Estrategias de análisis de datos: estadística descriptiva, análisis de contenido y triangulación
	\end{itemize}
}
{\footnotesize{\emph{Herramientas:} Moodle, Google Classroom, rúbricas analíticas, Mentimeter, Kahoot}}
\emptySeparator

\experience
{Enero 2024} {Analista de Datos Freelance — Datos Públicos e Indicadores Sociales}{Proyectos Personales}{Remoto}
{Enero 2024} {
	Desarrollo de proyectos de análisis de datos provenientes de fuentes estadísticas oficiales del Estado peruano para portafolio técnico.
	\begin{itemize}
		\item Dashboard interactivo de coyuntura económica regional con datos del INEI/BCRP/MEF usando Shiny
		\item Análisis exploratorio de encuestas nacionales (ENAHO, ENE) con R y Python
		\item Web scraping de datos económicos públicos y automatización de actualización de datasets oficiales
		\item Creación de visualizaciones interactivas para comunicación de hallazgos a audiencias no técnicas
		\item Documentación de proyectos en GitHub con README detallados y notebooks reproducibles
	\end{itemize}
}
{\footnotesize{\emph{Tecnologías:} R (dplyr, ggplot2, tidyr), Python (Pandas, Matplotlib, Seaborn), SQL, Git, Shiny, Streamlit}}

\end{experiences}
